\section{Trees cannot create scalar phase saving}
\label{sec:trees}

\begin{theorem}[Arbitrary phases on a tree]\label{thm:tree}
Suppose \(G_{|H|}\) is a tree and every tree edge is nonzero.  Then a
diagonal unitary \(D\) satisfies
\[
 D^*HD=|H|.
\]
Moreover the tree is bipartite, so another diagonal unitary switches
\(H\) to \(-|H|\).  Consequently
\begin{equation}\label{eq:tree-spectrum}
 \lambda_{\max}(H)=\rho(|H|),
 \qquad
 \lambda_{\min}(H)=-\rho(|H|),
 \qquad
 \|H\|=\rho(|H|).
\end{equation}
\end{theorem}

\begin{proof}
There are no cycle constraints, so the root propagation in
\cref{thm:holonomy} switches all phases to \(+1\).  Let
\(S=\diag(s_i)\), where \(s_i=+1\) on one bipartition class and
\(-1\) on the other.  Then \(S|H|S=-|H|\).  Hence \(H\) is unitarily
similar to both signs of \(|H|\), and \eqref{eq:tree-spectrum} follows.
\end{proof}

\begin{corollary}[Tree stop rule]\label{cor:tree-stop}
If every Perron-dominant connected component of the literal row graph
\(G_{|H|}\) is a tree, then the phase-frustration gap
\(\rho(|H|)-\|H\|\) is zero.  Any strict gap relative to a larger
majorant \(P\) is entirely amplitude/completion slack.
\end{corollary}

\begin{remark}[Not the TPC-78 tree]
TPC-78 uses a bipartite graph whose vertices are analysis coordinates
and completion variables.  The present tree has physical Hermitian rows
as vertices.  One can be a forest while the other contains many
cycles.  No tree conclusion transfers without constructing the actual
row matrix and its support.
\end{remark}
